% Zumdahl Chapter 8 free-response set -- 2 long (10) + 3 short (4) = 32 pts.
\documentclass{apchem-exam}

\apcunitword{Chapter}
\apcunit{8}
\apcunittitle{Bonding: General Concepts}
\apcdoctitle{Free-Response Set}
\apcsubtitle{Zumdahl \S8.1--\S8.13 \textbullet\ 5 questions \textbullet\ 32 points \textbullet\ 50 minutes}

\begin{document}
\apctitleblock

\begin{apcnote}[Scope]
A heavily examined chapter: \textbf{CED 2.1} (\S8.1--8.3), \textbf{2.2}
(\S8.5, \S8.8), \textbf{2.5} (Lewis diagrams, \S8.9--8.11), \textbf{2.6}
(resonance and formal charge, \S8.12), \textbf{2.7} (VSEPR, \S8.13),
\textbf{1.8} (\S8.4) and \textbf{6.7} (bond enthalpies, \S8.8).

Throughout, explanations of trends must be \textbf{Coulombic} --- charge,
distance, shielding. Rubrics here award zero for answers that rest on
``octet'' or ``stability'' language alone.
\end{apcnote}

\examsection{II}{Free Response}{5 questions \textbullet\ 32 points \textbullet\ 50 minutes}{\calculatorok}

% =====================================================================
\frq{long}{10}{Zumdahl \S8.10--\S8.13 \textbullet\ CED 2.5, 2.6, 2.7 \textbullet\ Lewis structures, resonance and shape}

Consider the nitrate ion, \ce{NO3-}.

\begin{parts}
  \item Draw a valid Lewis structure for \ce{NO3-}, showing all lone pairs
        and the overall charge. \workspace[3.4cm]
        \keyonly{\begin{apcanswer}Valence electrons:
        $5 + 3(6) + 1 = \mathbf{24}$ (the $+1$ is for the negative
        charge).

        Nitrogen is central. Draw one \ce{N=O} double bond and two
        \ce{N-O} single bonds. Each singly bonded oxygen carries three
        lone pairs; the doubly bonded oxygen carries two. Nitrogen has no
        lone pair. Enclose the structure in brackets with a $-$ charge
        outside.\end{apcanswer}}
  \item Explain why \ce{NO3-} is described by three resonance structures,
        and state what the true structure is.
        \answerlines[3]{The double bond can be placed on any of the three
        oxygens, giving three equivalent Lewis structures. None is correct
        on its own: the real ion is the average of all three, with the
        extra bonding electron pair delocalized over all three \ce{N-O}
        positions. The ion does not flip between them.}
  \item State and justify the prediction for the three \ce{N-O} bond
        lengths. \answerlines[3]{All three are \textbf{identical}, and
        each is intermediate between a normal \ce{N-O} single bond and a
        normal \ce{N=O} double bond. Delocalization spreads the extra pair
        equally, so every position has the same bond order, about
        $1\tfrac{1}{3}$.}
  \item Predict the shape and the \ce{O-N-O} bond angle, and name the
        model used. \answerlines[2]{\textbf{Trigonal planar},
        \SI{120}{\degree}, from the VSEPR model: three electron domains
        around nitrogen and no lone pair, so the domains spread as far
        apart as possible in a plane.}
\end{parts}

\begin{rubric}
  \rpoint{3}{(a) 1 pt 24 valence electrons including the charge; 1 pt one
    double and two single bonds; 1 pt all lone pairs shown with brackets
    and charge. Forgetting the extra electron for the charge caps this at
    1}
  \rpoint{3}{(b) 1 pt three equivalent placements of the double bond;
    1 pt the true structure is an average, not an alternation; 1 pt the
    word delocalized. Saying the ion ``resonates back and forth'' earns
    \textbf{0} for the second point}
  \rpoint{2}{(c) 1 pt all three equal; 1 pt intermediate between single
    and double, justified by delocalization}
  \rpoint{2}{(d) 1 pt trigonal planar with \SI{120}{\degree}; 1 pt three
    domains with no lone pair on the central atom}
\end{rubric}

% =====================================================================
\frq{long}{10}{Zumdahl \S8.5, \S8.8 \textbullet\ CED 2.2, 6.7 \textbullet\ bond enthalpy and lattice energy}

\begin{parts}
  \item Use the average bond enthalpies below to estimate $\Delta H$ for
        \[ \ce{CH4(g) + 2O2(g) -> CO2(g) + 2H2O(g)} \]
        \ce{C-H} 413, \ce{O=O} 495, \ce{C=O} 799, \ce{O-H} 467
        \si{\kilo\joule\per\mole}. \workspace[3.6cm]
        \keyonly{\begin{apcanswer}Bonds \textbf{broken} (energy in):
        $4(413) + 2(495) = 1652 + 990 = \num{2642}$~kJ

        Bonds \textbf{formed} (energy out):
        $2(799) + 4(467) = 1598 + 1868 = \num{3466}$~kJ

        $\Delta H = \text{broken} - \text{formed} = 2642 - 3466 =
        \mathbf{-824}$~kJ

        Negative, so exothermic --- more energy is released forming the
        new bonds than is required to break the old
        ones.\end{apcanswer}}
  \item Explain why this method gives only an estimate.
        \answerlines[3]{The values are \emph{average} bond enthalpies
        taken across many different compounds. The actual strength of a
        given \ce{C-H} bond depends on its molecular environment, so the
        computed total differs from the true value. The method also
        assumes all species are gaseous.}
  \item Lithium fluoride and potassium fluoride are both ionic. Predict
        which has the larger lattice energy and justify in Coulombic
        terms. \workspace[3cm]
        \keyonly{\begin{apcanswer}\textbf{\ce{LiF}}. Both compounds pair
        singly charged ions, so the charge product is the same and cannot
        be the deciding factor.

        \ce{Li+} is a much smaller cation than \ce{K+}, so in \ce{LiF} the
        ions approach more closely and the internuclear distance $d$ is
        shorter. Coulombic attraction varies as $1/d$, so the shorter
        distance gives the stronger attraction and the larger lattice
        energy.\end{apcanswer}}
  \item State the effect on lattice energy of replacing \ce{F-} with
        \ce{O^2-} in a compound of the same cation, and give the reason.
        \answerlines[2]{The lattice energy increases substantially. The
        charge product doubles from $(+1)(-1)$ to $(+1)(-2)$, and
        Coulombic attraction is proportional to the product of the
        charges.}
\end{parts}

\begin{rubric}
  \rpoint{4}{(a) 1 pt bonds broken \num{2642}; 1 pt bonds formed
    \num{3466}; 1 pt broken minus formed in the correct sense; 1 pt
    $-824$~kJ. Reversing the subtraction gives $+824$ and earns at most 2}
  \rpoint{2}{(b) 1 pt the values are averages over many compounds; 1 pt
    actual bond strength depends on environment}
  \rpoint{3}{(c) 1 pt \ce{LiF}; 1 pt noting the charges are equal so
    distance decides; 1 pt smaller cation gives shorter $d$ and stronger
    attraction. ``\ce{LiF} is more stable'' earns \textbf{0}}
  \rpoint{1}{(d) increases, because the charge product doubles}
\end{rubric}

% =====================================================================
\frq{short}{4}{Zumdahl \S8.2--\S8.3 \textbullet\ CED 2.1 \textbullet\ bond polarity and molecular polarity}

\begin{parts}
  \item \ce{CO2} and \ce{H2O} both contain polar bonds, yet only one
        molecule is polar. Identify which, and explain.
        \workspace[3.2cm]
        \keyonly{\begin{apcanswer}\textbf{Water} is polar; carbon dioxide
        is not.

        In \ce{CO2} the molecule is \textbf{linear}, so the two equal
        \ce{C=O} bond dipoles point in exactly opposite directions and
        cancel, leaving no net dipole.

        In \ce{H2O} the two lone pairs on oxygen force a \textbf{bent}
        shape. The two \ce{O-H} bond dipoles no longer oppose each other,
        so they add to give a net dipole pointing toward the
        oxygen.\end{apcanswer}}
  \item State what determines whether a \emph{bond} is polar.
        \answerlines[2]{A difference in electronegativity between the two
        bonded atoms: the more electronegative atom draws the shared pair
        closer, giving it a partial negative charge.}
\end{parts}

\begin{rubric}
  \rpoint{3}{(a) 1 pt water identified; 1 pt \ce{CO2} linear with dipoles
    cancelling; 1 pt \ce{H2O} bent so dipoles do not cancel. An answer
    that never mentions shape earns at most 1}
  \rpoint{1}{(b) electronegativity difference}
\end{rubric}

% =====================================================================
\frq{short}{4}{Zumdahl \S8.13 \textbullet\ CED 2.7 \textbullet\ VSEPR}

For each species, state the number of electron domains on the central
atom, the molecular shape, and the approximate bond angle.

\begin{parts}
  \item \ce{NH3} \answerlines[2]{Four domains (three bonding, one lone
        pair); \textbf{trigonal pyramidal}; slightly less than
        \SI{109.5}{\degree}, about \SI{107}{\degree}.}
  \item \ce{SO2} \answerlines[2]{Three domains (two bonding, one lone
        pair); \textbf{bent}; slightly less than \SI{120}{\degree}.}
  \item Explain why the bond angle in \ce{NH3} is smaller than the ideal
        tetrahedral angle. \answerlines[2]{A lone pair is held by only one
        nucleus, so it spreads out closer to the central atom and occupies
        more angular space than a bonding pair. It presses the three
        \ce{N-H} bonds together, reducing the angle between them.}
\end{parts}

\begin{rubric}
  \rpoint{1}{(a) four domains, trigonal pyramidal, angle below
    \SI{109.5}{\degree}}
  \rpoint{1}{(b) three domains, bent, angle below \SI{120}{\degree}}
  \rpoint{2}{(c) 1 pt a lone pair occupies more space than a bonding pair;
    1 pt it compresses the bonding angles. Note the shape is named from
    the \emph{atoms}, not the domains}
\end{rubric}

% =====================================================================
\frq{short}{4}{Zumdahl \S8.4 \textbullet\ CED 1.8 \textbullet\ isoelectronic ions}

The ions \ce{N^3-}, \ce{O^2-}, \ce{F-}, \ce{Na+} and \ce{Mg^2+} all have
ten electrons.

\begin{parts}
  \item Arrange them in order of increasing ionic radius.
        \answerlines[2]{\ce{Mg^2+} $<$ \ce{Na+} $<$ \ce{F-} $<$ \ce{O^2-}
        $<$ \ce{N^3-}}
  \item Justify the order. \workspace[2.8cm]
        \keyonly{\begin{apcanswer}Every one of these ions has the same ten
        electrons, so shielding is essentially identical across the whole
        set. The only quantity that changes is the \textbf{nuclear
        charge}.

        \ce{Mg^2+} has 12 protons pulling on those ten electrons;
        \ce{N^3-} has only 7. The greater the proton count, the more
        strongly the same electron cloud is drawn in, and the smaller the
        ion. Radius therefore falls steadily as nuclear charge
        rises.\end{apcanswer}}
\end{parts}

\begin{rubric}
  \rpoint{2}{(a) the full order; 1 pt if only one pair is transposed}
  \rpoint{2}{(b) 1 pt same number of electrons so shielding is constant;
    1 pt increasing nuclear charge pulls the cloud in more tightly.
    ``Anions are bigger than cations'' with no Coulombic reasoning earns
    \textbf{0}}
\end{rubric}

\examtotal

\end{document}
